\chapter{Gestaltung und Formatregeln}\label{kap:design}

Die Gestaltung ist zurückhaltend und zeitungsnah. Die Daten haben Vorrang,
und Farbe trägt Bedeutung, statt zu schmücken. Alle Regeln stehen in einer
handgeschriebenen Stilvorlage, \datei{RCode/site_assets/site.css}. Sie
kommt ohne Framework aus.

\section{Schriften}

Alle drei Schriften liegen als WOFF2 im Projekt, nur mit lateinischem
Zeichensatz, und werden selbst ausgeliefert. Es gibt keine Anfragen an
Schriftendienste.

\begin{center}
\small
\begin{tabularx}{\textwidth}{@{}llX@{}}
\toprule
Rolle & Schrift & Einsatz \\
\midrule
Serif & Source Serif 4 (variabel, 200--900) & Fließtext, Überschriften, Wortmarke, Vereinsnamen; Mediävalziffern im Text \\
Sans  & Inter Tight (400--700) & Navigation, Kapitälchen-Überschriften, Tabellenköpfe, Legenden, Fuß \\
Mono  & JetBrains Mono (400--700) & alle Zahlen in Tabellen, Heatmap, Ergebnisse, ELO (Tabellenziffern) \\
\bottomrule
\end{tabularx}
\end{center}

Dieses Handbuch ist in Source Serif Pro, Inter und Source Code Pro gesetzt,
den nächsten Verwandten in \TeX\ Live.

\section{Farben}

\newcommand{\farbe}[3]{\textcolor{#1}{\rule{1.6em}{0.9em}} & \code{#2} & \code{#3}}
\begin{center}
\small
\begin{tabularx}{\textwidth}{@{}cllX@{}}
\toprule
& Token & Wert & Verwendung \\
\midrule
\farbe{ink}{--ink}{\#15130F} & Haupttext \\
\farbe{ink2}{--ink-2}{\#312E29} & Sekundärtext, Navigation \\
\farbe{ink3}{--ink-3}{\#6B655C} & Legenden, Fuß, Kapitälchen \\
\farbe{rule}{--rule}{\#D4CEBF} & Linien, Rahmen, Linien im Verlauf \\
\farbe{rulesoft}{--rule-soft}{\#E3DED0} & feine Linien, Remis-Segment \\
\farbe{surface2}{--surface-2}{\#F5F3EE} & Hinweisflächen, Winterpause, gewählte Zeile \\
\farbe{akzent}{--akzent}{\#B0261E} & Rot: Unterstreichung, Strich unter der Wortmarke, Abstieg, negative Änderung, Fokusrahmen \\
\farbe{tinte}{--tinte}{\#1F3A4D} & Tintenblau: Endfarbe der Heatmap, Heimsieg, hervorgehobene Linie, positive Änderung \\
\farbe{ocker}{--ocker}{\#B8893A} & Auswärtssieg, Relegation, Hinweisbalken \\
\farbe{gruen}{--gruen}{\#3F7D3A} & Aufstiegszone \\
\bottomrule
\end{tabularx}
\end{center}

Die Seite hat bewusst nur ein helles Erscheinungsbild. Einen Dunkelmodus
gibt es nicht.

\section{Satz und Raster}

\begin{itemize}
  \item Fließtext: Serif, 17~px, Zeilenhöhe 1{,}58. Der Inhalt ist höchstens
        1\,080~px breit, Einleitungstexte höchstens 64 Zeichen.
  \item Abschnitte: 56~px Abstand und eine feine Linie. Jede Sektion beginnt
        mit einer kleinen Überschrift in Kapitälchen (\enquote{PROGNOSE},
        \enquote{TABELLE}, \dots), danach folgt die Überschrift in Serif.
  \item Links: dunkle Schrift mit roter Unterstreichung.
  \item Umbrüche bei \textbf{760~px} (kleinere Abstände, Spielzeilen
        zweispaltig, Verlaufsdiagramm scrollbar) und \textbf{520~px}
        (Ligatabelle ohne \enquote{Sp.} und \enquote{Tordiff.}, Heatmap mit
        fester Spaltenbreite, Werte unter 1~Prozent nur als Farbe).
\end{itemize}

\section{Zahlen und Prozente}\label{sec:prozent}

Prozentwerte formatiert eine einzige Funktion, \code{prozent()}
(\datei{RCode/render_helpers.R:32--47}):

\begin{center}
\small
\begin{tabular}{@{}ll@{}}
\toprule
Wahrscheinlichkeit $x$ & Anzeige \\
\midrule
$x = 0$ & \enquote{0} (in der Heatmap: leere Zelle) \\
$0 < x < 0{,}01$ & \enquote{<1} \\
$0{,}01 \le x \le 0{,}99$ & $\operatorname{round}(100x)$ \\
$0{,}99 < x < 1$ & \enquote{>99} \\
$x = 1$ & \enquote{\checkmark} \\
\bottomrule
\end{tabular}
\end{center}

Damit steht nie \enquote{100} oder \enquote{0} für etwas, das in einzelnen
Simulationen doch anders ausging. Weitere Regeln:

\begin{itemize}
  \item Dezimalkomma überall. Negative Zahlen tragen das echte Minuszeichen
        (U+2212), null trägt \enquote{$\pm$}.
  \item Wo ein Prozentzeichen steht, trennt es ein geschütztes Leerzeichen von
        der Zahl (\enquote{12\,\%}). Tabellen verzichten auf das Zeichen und
        sagen es in der Legende.
  \item Die Ergebnis-Matrix zeigt eine Nachkommastelle, alle anderen
        Prozentangaben sind ganzzahlig.
  \item Zeiten stehen in Berliner Zeit mit \enquote{Uhr}, im Fuß mit
        \enquote{MEZ} oder \enquote{MESZ}.
\end{itemize}

\section{Barrierefreiheit}

\begin{itemize}
  \item Die Heatmap trägt \code{aria-label}, und die Balken sind mit
        \code{role="img"} und einer gesprochenen Beschreibung versehen
        (\enquote{Sieg Heim 62\,\%, Remis 22\,\%, Sieg Gast 16\,\%}).
  \item Kürzel sind als \code{<abbr>} mit dem vollen Namen ausgezeichnet und
        per Tastatur und Antippen erreichbar.
  \item Das Verlaufsdiagramm hat eine Textbeschreibung, und die Tabelle
        darunter enthält dieselben Zahlen als zugängliche Alternative.
  \item Alle bedienbaren Elemente haben einen sichtbaren roten Fokusrahmen.
        Die einzige Animation (das Aufklappen der Ergebnis-Matrix) entfällt
        bei \code{prefers-reduced-motion}.
\end{itemize}
